\section{Постановка задачи}

\textbf{Цель работы:} разработать библиотеку для автоматической генерации мапперов на основе аннотаций с использованием KSP, реализующую метапрограммирование на этапе компиляции, которая позволит сократить объём шаблонного кода за счёт создания вспомогательных функций, дополняющих исходные мапперы недостающими параметрами преобразования.

Для достижения поставленной цели необходимо решить следующие задачи:

\begin{enumerate}
    \item Изучить документацию и API библиотеки KSP \cite{refg, refj}, освоить механизмы обхода символов и генерации файлов как основу метапрограммирования.
    \item Изучить библиотеку KotlinPoet \cite{refk} для программного создания исходного кода на Kotlin.
    \item Разработать аннотацию \texttt{@CastCastleMapper} для того, чтобы обеспечить обработку помеченных ею частей кода
    \item Реализовать KSP-процессор, включающий поиск аннотированных классов и функций, анализ их структуры.
    \item Реализовать компонент для разрешения зависимостей между мапперами и построения моделей данных.
    \item Реализовать генератор кода, создающий вспомогательные функции.
    \item Добавить поддержку вложенных объектов и коллекций с рекурсивным вызовом мапперов.
    \item Написать интеграционные тесты и провести ручное тестирование.
\end{enumerate}
